% ============================================================================
% Fichier  : partie1/P01_C02_perception.tex
% Titre    : Percevoir le Risque : Biais et Culture
% Auteur   : MINKA MI NGUIDJOÏ Thierry Emmanuel
% Version  : 1.0 -- novembre 2026
% ============================================================================
% Sources : Tversky et Kahneman (1974) ; Kahneman (2011) ; Slovic (1987) ;
% Janis (1972) ; Vaughan (1996) ; Reason (1997) ; Hubbard et Seiersen (2016).
% ============================================================================

\chapter{Percevoir le Risque : Biais et Culture}
\label{chap:perception}
\chaptermark{Percevoir le risque}

\epigraph{%
    « Ce qui nous attire des ennuis, ce n'est pas ce que nous ignorons,\\
    c'est ce que nous tenons pour certain et qui ne l'est pas. »%
}{Formule d'usage dans la gestion des risques}

\niveauChapitre{Fondamental}{Chapitre~\ref{chap:risque}.}

\begin{boxobjectifs}
\begin{itemize}
    \item Reconnaître les heuristiques de disponibilité, de représentativité
          et d'ancrage dans un jugement sur le risque.
    \item Expliquer pourquoi le public et les experts ne hiérarchisent pas
          les risques de la même manière.
    \item Identifier les biais collectifs~: pensée de groupe, normalisation
          de la déviance.
    \item Décrire une culture juste et son rôle dans la remontée des
          signaux.
    \item Calibrer ses propres estimations par des intervalles.
\end{itemize}
\end{boxobjectifs}

\bigskip

% ============================================================================
\section*{Question directrice : pourquoi « nous n'avons jamais été attaqués » a-t-il convaincu~?}
% ============================================================================

En février, l'argument a suffi à reporter le second facteur. En avril,
après l'attaque, le comité de direction a voté sans débat un budget de
sécurité quatre fois supérieur à celui de l'année précédente. Dans les
deux cas, la décision a suivi la mémoire récente plus que l'analyse. Ce
chapitre explique pourquoi, et comment s'en protéger.

\medskip
\textbf{Plan du chapitre.}
\begin{enumerate}
    \item Deux manières de penser (section~\ref{sec:c02-systemes}).
    \item Les heuristiques (section~\ref{sec:c02-heuristiques}).
    \item La perception du risque (section~\ref{sec:c02-perception}).
    \item Les biais collectifs (section~\ref{sec:c02-collectifs}).
    \item La culture du risque (section~\ref{sec:c02-culture}).
    \item Se calibrer (section~\ref{sec:c02-calibrer}).
\end{enumerate}

\separateur

% ============================================================================
\section{Deux Manières de Penser}
\label{sec:c02-systemes}
% ============================================================================

Daniel Kahneman décrit deux modes de pensée~: un mode rapide, intuitif,
automatique, et un mode lent, délibéré, coûteux en
effort~\cite{kahneman2011}. Le premier nous permet de réagir vite~; il
répond souvent à une question plus facile que celle qui est posée. À la
question « quelle est la probabilité d'une attaque cette année~? », il
répond « me souviens-je d'une attaque récente~? ».

\begin{boxdef}
\termecle{Heuristique}~: règle de jugement simple, le plus souvent
inconsciente, qui donne vite une réponse acceptable, mais produit des
erreurs systématiques, les \termecle{biais}.
\end{boxdef}

% ============================================================================
\section{Les Heuristiques}
\label{sec:c02-heuristiques}
% ============================================================================

Amos Tversky et Daniel Kahneman ont décrit en 1974 trois heuristiques du
jugement sous incertitude~\cite{tversky1974}.

\begin{tableaucours}{L{3.0cm}L{4.6cm}L{4.6cm}}{%
    \tableauHeader{Heuristique} &
    \tableauHeader{Principe} &
    \tableauHeader{Chez LiZbiZ}
}
Disponibilité & On juge probable ce dont on se souvient facilement &
« Jamais attaqués » en février~; « tout est menacé » en avril \\
Représentativité & On juge sur la ressemblance avec un stéréotype, en
négligeant la fréquence de base & « Un rançongiciel vise les banques, pas
la logistique » \\
Ancrage & Une valeur initiale, même arbitraire, tire l'estimation &
Le budget de l'an passé sert de point de départ au suivant \\
\end{tableaucours}
\vspace{-0.8em}
\begin{center}{\small\itshape Tableau~2.1. Les trois heuristiques de Tversky et Kahneman}\end{center}

\begin{boxalert}
\textbf{Violation typique.} Estimer une fréquence d'attaque par sa propre
expérience (« jamais vu ici »).

\smallskip
\textbf{Règle~:} partir d'une \emph{fréquence de base} externe (études
publiées, données du secteur, signalements du centre de réponse national),
puis l'ajuster par ce qui distingue l'organisation. L'absence d'attaque
constatée peut aussi signifier une absence de détection.
\end{boxalert}

S'y ajoutent l'\termecle{excès d'optimisme} (croire que les malheurs
arrivent aux autres), le \termecle{biais de normalité} (penser qu'une
situation inhabituelle va revenir d'elle-même à la normale, comme le
technicien qui redémarre les serveurs le 5~mars au soir) et le
\termecle{biais rétrospectif} (juger après coup qu'un événement était
évident).

% ============================================================================
\section{La Perception du Risque}
\label{sec:c02-perception}
% ============================================================================

Paul Slovic a montré que le public ne classe pas les risques selon leur
fréquence de décès, mais selon deux dimensions~: le caractère
\emph{redouté} (catastrophique, involontaire, incontrôlable) et le
caractère \emph{inconnu} (nouveau, invisible, aux effets
différés)~\cite{slovic1987}. Un risque rare mais redouté et inconnu
(une cyberattaque massive) mobilise plus qu'un risque fréquent et
familier (les pannes).

\begin{boxscenario}{Opération GAROUA. Deux budgets}
\textbf{Constatation.} Budget de sécurité de 2025~: 45~millions de FCFA~;
budget voté en avril 2026~: 180~millions, sans analyse des risques
préalable.

\medskip
\textbf{Analyse} (confiance moyenne). Le premier budget était ancré sur
l'historique et l'absence d'incident~; le second, sur l'effroi d'une
crise récente. Ni l'un ni l'autre n'a été dérivé d'une analyse. Le risque
est que les 180~millions financent des mesures visibles plutôt que les
plus efficaces~: le registre du Chapitre~\ref{chap:traiter} doit les
réaffecter.
\end{boxscenario}

% ============================================================================
\section{Les Biais Collectifs}
\label{sec:c02-collectifs}
% ============================================================================

Les groupes ne corrigent pas toujours les biais individuels~; ils peuvent
les amplifier.

\begin{tableaucours}{L{3.6cm}L{8.6cm}}{%
    \tableauHeader{Biais} &
    \tableauHeader{Description et parade}
}
Pensée de groupe & Un groupe soudé recherche l'unanimité au détriment de
l'examen des options (Janis). Parade~: désigner un contradicteur, recueillir
les avis par écrit avant la discussion~\cite{janis1972} \\
Normalisation de la déviance & Un écart toléré sans conséquence devient la
norme (Vaughan, à propos de Challenger). Parade~: tracer les dérogations et
les revoir~\cite{vaughan1996} \\
Effet d'autorité & L'avis du plus haut placé clôt le débat. Parade~: faire
parler le plus haut placé en dernier \\
\end{tableaucours}
\vspace{-0.8em}
\begin{center}{\small\itshape Tableau~2.2. Trois biais collectifs}\end{center}

\textbf{Exemple.} Le compte partagé du prestataire sans second facteur
était connu depuis septembre~; il avait été accepté « provisoirement »,
puis plus personne n'en parlait. C'est une normalisation de la déviance~:
une exception devenue habitude.

% ============================================================================
\section{La Culture du Risque}
\label{sec:c02-culture}
% ============================================================================

James Reason décrit une \termecle{culture juste}~: une organisation où l'on
signale les erreurs et les quasi-accidents sans crainte de sanction
injuste, tout en sanctionnant les comportements délibérément dangereux~;
elle conditionne la remontée des signaux faibles~\cite{reason1997}.

\begin{boxmetier}{Gestionnaire de risques}
\textbf{Rendre le signalement facile et utile.} Le gestionnaire ouvre un
canal simple pour signaler un quasi-incident, répond à chaque signalement,
et publie ce qui a changé grâce aux signalements. Un salarié qui ne voit
jamais d'effet cesse de signaler.
\end{boxmetier}

% ============================================================================
\section{Se Calibrer}
\label{sec:c02-calibrer}
% ============================================================================

Douglas Hubbard et Richard Seiersen montrent qu'un expert peut apprendre à
donner des estimations \termecle{calibrées}~: des intervalles à 90~\% qui
contiennent la vraie valeur environ neuf fois sur dix~\cite{hubbard2016cyber}.
Sans entraînement, la plupart des personnes donnent des intervalles trop
étroits~: elles sont trop sûres d'elles.

\begin{boxprocedure}
\textbf{Estimer par un intervalle calibré}
\begin{enumerate}
    \item Donner une borne basse et une borne haute telles que l'on soit
          sûr à 90~\% que la valeur est entre les deux.
    \item Se demander~: « parierais-je autant sur mon intervalle que sur
          une roue qui gagne neuf fois sur dix~? » Élargir si la roue
          semble préférable.
    \item Chercher une raison pour laquelle la valeur serait hors de
          l'intervalle, de chaque côté.
    \item Après coup, compter combien d'intervalles contenaient la valeur,
          et corriger sa largeur.
\end{enumerate}
\end{boxprocedure}

\separateur

\begin{boxresume}
\textbf{Ce qu'il faut retenir de ce chapitre~:}
\begin{itemize}
    \item \textbf{Le jugement rapide} substitue une question facile à la
          question posée.
    \item \textbf{Disponibilité, représentativité, ancrage} faussent
          l'estimation~; on part d'une fréquence de base.
    \item \textbf{Le public} juge sur l'effroi et l'inconnu~; l'expert doit
          le savoir pour communiquer.
    \item \textbf{Les groupes} amplifient les biais~; contradicteur et avis
          écrits les limitent.
    \item \textbf{Une culture juste} fait remonter les signaux~;
          \textbf{l'intervalle calibré} rend l'estimation honnête.
\end{itemize}
\end{boxresume}

\begin{boxinfo}
\textbf{Points de contrôle de l'élève.} Je sais~: nommer et reconnaître
trois heuristiques~; expliquer les dimensions de Slovic~; décrire la
pensée de groupe et la normalisation de la déviance~; définir une culture
juste~; donner un intervalle calibré.
\end{boxinfo}

\bigskip

\begin{boxexo}[title={\ding{45}\quad Exercice~2.1. Quelle heuristique~? \niveaubase}]
Nommez l'heuristique~: « depuis l'attaque d'une banque voisine, je trouve
les rançongiciels très probables »~; « le budget de cette année sera celui
de l'an dernier plus 10~\% ».
\end{boxexo}

\begin{boxexo}[title={\ding{45}\quad Exercice~2.2. Fréquence de base \niveaubase}]
Pourquoi l'absence d'attaque constatée chez LiZbiZ avant mars ne prouve-t-elle
pas l'absence d'attaque~?
\end{boxexo}

\begin{boxexo}[title={\ding{45}\quad Exercice~2.3. Biais de normalité \niveaubase}]
En quoi le redémarrage des serveurs le 5~mars à 22~h~05 illustre-t-il le
biais de normalité~?
\end{boxexo}

\begin{boxexo}[title={\ding{45}\quad Exercice~2.4. Slovic \niveaumoyen}]
Placez sur les deux axes de Slovic~: une panne d'électricité~; une fuite
de données personnelles~; une inondation d'entrepôt~; une fraude interne.
\end{boxexo}

\begin{boxexo}[title={\ding{45}\quad Exercice~2.5. Le comité \niveaumoyen}]
Proposez trois règles de fonctionnement du comité des risques de LiZbiZ
pour limiter la pensée de groupe et l'effet d'autorité.
\end{boxexo}

\begin{boxexo}[title={\ding{45}\quad Exercice~2.6. Dérogations \niveaumoyen}]
Concevez un registre des dérogations qui empêche la normalisation de la
déviance~: colonnes, propriétaire, durée, revue.
\end{boxexo}

\begin{boxexo}[title={\ding{45}\quad Exercice~2.7. Culture juste \niveaumoyen}]
Un agent a cliqué sur un lien d'hameçonnage et l'a signalé aussitôt.
Comment la direction doit-elle réagir dans une culture juste~? Et s'il
l'avait caché~?
\end{boxexo}

\begin{boxexo}[title={\ding{45}\quad Exercice~2.8. Calibrer \niveaumoyen}]
Donnez un intervalle à 90~\% pour~: la population de Douala~; l'année de
publication de l'ISO~31000 (première édition)~; la distance Douala-Garoua
par la route. Vérifiez ensuite.
\end{boxexo}

\begin{boxexo}[title={\ding{45}\quad Exercice~2.9. Le budget d'avril \niveauexpert}]
Proposez à M.~ETEME une méthode en quatre étapes pour affecter les
180~millions de FCFA sans céder à l'effroi.
\end{boxexo}

\begin{boxexo}[title={\ding{45}\quad Exercice~2.10. Un atelier \niveauexpert}]
Concevez un atelier d'une heure pour calibrer les dix experts qui
estimeront les risques de LiZbiZ~: questions, déroulé, mesure du
résultat.
\end{boxexo}

\bigskip

\begin{boxinfo}
\textbf{Pour aller plus loin.}
\begin{itemize}
    \item A.~Tversky et D.~Kahneman~\cite{tversky1974}~; D.~Kahneman~\cite{kahneman2011}.
    \item P.~Slovic~\cite{slovic1987}~; J.~Reason~\cite{reason1997}.
    \item D.~Hubbard et R.~Seiersen~\cite{hubbard2016cyber}.
\end{itemize}
\end{boxinfo}

\begin{boxinfo}
\textbf{Transition.} Pour que les décisions sur les risques ne dépendent
pas de l'humeur du moment, il faut des rôles et des règles. Le
Chapitre~\ref{chap:gouvernance} présente la gouvernance des risques.
\end{boxinfo}
